% Part: incompleteness
% Chapter: theories-computability
% Section: consis-with-q

\documentclass[../../../include/open-logic-section]{subfiles}

\begin{document}

\olfileid{inc}{tcp}{con}

\olsection{$\Th{Q}$తో అవైరుధ్యంగా ఉన్న సిద్ధాంతాలు అనిర్ణయనీయాలు}

కింది సిద్ధాంతం $\Th{Q}$ మాత్రమే అనిర్ణయనీయం కాదనీ, నిజానికి $\Th{Q}$తో
విభేదించని ప్రతి సిద్ధాంతమూ అనిర్ణయనీయమేననీ చెబుతుంది.

\begin{thm}
అంకగణిత భాషలోని సిద్ధాంతం~$\Th{T}$, $\Th{Q}$తో అవైరుధ్యంగా ఉందనుకుందాం
(అంటే $\Th{T} \cup \Th{Q}$ అవైరుధ్యమైనది). అప్పుడు $\Th{T}$
అనిర్ణయనీయం.
\end{thm}

\begin{proof}
$\Th{Q}$కు $!Q_1$, \dots, $!Q_8$ అనే పరిమిత స్వీకృతాల సమితి ఉందని
గుర్తుంచుకోండి. వీటి స్థానంలో $!E = !Q_1 \land \dots \land !Q_8$ అనే ఒకే
స్వీకృతాన్ని కూడా పెట్టవచ్చు.

$\Th{T}$ అనేది~$\Th{Q}$తో అవైరుధ్యంగా ఉన్న నిర్ణయించదగిన సిద్ధాంతం అని
అనుకుందాం. కిందిదాన్ని తీసుకుందాం:
\[
C = \Setabs{!A}{\Th{T} \Proves !E \lif !A}.
\]
$C$ అనేది $\Th{Q}$, $\Th{\bar Q}$లను గణనీయంగా వేరుపరిచే సమితి అవుతుందని
చూపి వైరుధ్యాన్ని పొందుతాం. ముందుగా, $!A$, $\Th{Q}$లో ఉంటే, $\Th{Q}$
స్వీకృతాల నుంచి $!A$ నిరూపించదగినది; నిగమన సిద్ధాంతం వల్ల మొదటిస్థాయి
తర్కంలో $!E \lif !A$కు \tetoken{ఒక వ్యుత్పత్తి}{!!a{derivation}} ఉంటుంది. కాబట్టి $!A$,~$C$లో ఉంటుంది.

మరోవైపు, $!A$, $\Th{\bar Q}$లో ఉంటే, మొదటిస్థాయి తర్కంలో $!E \lif
\lnot !A$కు ఒక నిరూపణ ఉంటుంది. $\Th{T}$ కూడా $!E \lif !A$ను నిరూపిస్తే,
$\Th{T}$ అనేది $\lnot !E$ను నిరూపిస్తుంది; అప్పుడు $\Th{T} \cup \Th{Q}$
వైరుధ్యమైనదవుతుంది. కానీ $\Th{T} \cup \Th{Q}$ అవైరుధ్యమైనదని ఊహిస్తున్నాం.
కాబట్టి $\Th{T}$ అనేది $!E \lif !A$ను నిరూపించదు; అందువల్ల $!A$,~$C$లో
ఉండదు.

$!A$, $\Th{Q}$లో ఉంటే $C$లో ఉంటుందనీ, $!A$, $\Th{\bar Q}$లో ఉంటే
$\Complement{C}$లో ఉంటుందనీ చూపాం. కాబట్టి $C$ ఒక గణనీయ వేరుపరిచే సమితి;
ఇదే మనం వెతుకుతున్న వైరుధ్యం.
\end{proof}

ఈ సిద్ధాంతం చాలా శక్తిమంతమైనది. ఉదాహరణకు, దీని నుంచి కిందిది అనుగమిస్తుంది:

\begin{cor}
  అంకగణిత భాషకు మొదటిస్థాయి తర్కం (అంటే
  $\Setabs{!A}{\text{$!A$ మొదటిస్థాయి తర్కంలో నిరూపించదగినది}}$ అనే సమితి)
  అనిర్ణయనీయం.
\end{cor}

\begin{proof}
మొదటిస్థాయి తర్కం ఖాళీ సమితి~$\emptyset$ యొక్క పర్యవసానాల సమితి;
అది~$\Th{Q}$తో అవైరుధ్యంగా ఉంటుంది.
\end{proof}

\end{document}
